\scanpage{131}
% archive foliation: 130
Dans les éléments de structures canoniques associés à la courbe analytique $X$, nous allons \unsure{retenir} maintenant ceux qui ``gardent un sens algébrique'' lorsque $X$ ``est'' algébrique --- tt au moins dans le cas propre, grâce à \unsure{Gaga|GAGA}.
Ce sont

a) Fibré en droites projectives $\inserted{P_X}$ sur $X$

b) connexion (plate) $c_X$ sur $P_X$ (\unsure{rel.} $X/\mathbb C$)

c) section $\sigma$ de $P_X$

[d/ Un isom. de schémas formels sur $X$
\[\widehat P_X\xrightarrow[\sim]{\widehat p_X}\widehat X_X\]
(complétés formels le long des sections marquées $\sigma$ et $\Delta_X$) respectant les connexions.

D'ailleurs, l'élément de structure d) est contenu dans b) et \unsure{d)|c)}, moyennant la condition (\unsure{équivaut à} \deleted{\unsure{une liaison algébrique}}\inserted{$\exists$ iso d) compatible} avec conn) :

$1^\circ)$ \uline{La section $\sigma$ est nulle part stationnaire} (pour la connexion donnée $c$ sur $P_X$\ldots)]

Même si $X$ n'est pas propre, \unsure{le topo} de Deligne montre l'existence et de l'unicité de structures algébriques a) b) \inserted{(mais pas l'algéb. de c !)} déduites des constructions transcendantes du th. 2, lorsqu'on \deleted{fait}\inserted{pose} de plus la condition :

\deleted{\unsure{2}} La connexion $c_X$ de $P_X$ est \uline{régulière} à l'infini.

Mais nous allons étudier l'existence \unsure{tout} établir l'unicité, et préciser le Pb d'existence --- par une autre voie.
